
This work demonstrates that Model Context Protocol traces encode researcher reasoning in natural language, enabling automated semantic validation without manual test specification. Three contributions emerge from empirical evaluation on 20 research pipeline executions containing 596 tool calls.

First, MCP traces are rich semantic artifacts suitable for zero-annotation analysis. The 97.48\% trace completeness and 97.48\% natural language presence rates validate the foundational assumption that MCP logging captures sufficient reasoning content for semantic extraction. This extends prior work on MCP tool composition by demonstrating that traces themselves have value beyond runtime orchestration---they encode the reasoning behind tool calls.

Second, zero-annotation semantic extraction feasibility is established through Layers 1-2. LLM-based extraction achieved 82.7\% recall and 86.3\% precision with substantial inter-rater agreement (Cohen's $\kappa$=0.716) when validated against independent human annotation, using only zero-shot prompting and multi-vote consensus. This complements recent work on reducing annotation cost by achieving zero-annotation via trace analysis. The validated components---syntactic validation and semantic extraction---can be deployed immediately for trace quality monitoring and assumption archaeology from execution logs.

Third, a critical limitation in constraint inference is identified with a clear path forward. Layer 3 (semantic similarity matching) failed to detect assumption-evidence contradictions, achieving 0\% recall across 1,200 pairs. This negative result establishes a boundary condition: cosine similarity on sentence embeddings optimizes for topic relatedness (paraphrase detection), not logical contradiction. Contradictory statements like ''effective rank will decrease'' and ''effective rank increased by 6.02\%'' share high semantic similarity because they discuss the same entity and concept despite opposite polarity. This finding aligns with natural language inference literature and points toward entailment models (BERT-NLI, DeBERTa-MNLI) or LLM-based contradiction detection as necessary alternatives.

\textbf{Limitations and scoping.} Results are validated on research pipelines where natural language content is prevalent. Production MCP use cases with structured data processing may exhibit lower NL content, requiring domain-specific validation. The constraint inference validation used limited ground truth (one documented failure case), though threshold tuning across 1,200 pairs suggests systematic method mismatch rather than data insufficiency. The full three-layer framework's predicted failure detection rate remains unverified pending Layer 3 redesign. The refined claim acknowledges this: two layers validated, one requires methodological improvement.

\textbf{Future work.} Immediate extensions include replacing semantic similarity with entailment models (BERT-NLI, DeBERTa-MNLI) for contradiction detection, completing end-to-end validation on the 20-trace dataset, and testing generalization to non-research MCP pipelines. Longer-term directions include cross-project constraint learning by training entailment models on assumption-evidence-contradiction triplets from multiple pipelines, real-time validation by streaming MCP traces with incremental extraction, and research archaeology that infers unstated hypotheses and assumptions from historical pipeline execution logs.

By treating MCP traces as semantic artifacts, this work opens a validation paradigm where researcher reasoning becomes programmatically verifiable. The path forward is clear: two layers validated with principled limitations, one requires redesign with established alternatives, and a complete framework awaits integration.
